\lysh{15038_SOLUTION}{1}
\text{ΛΥΣΗ}

\begin{enumerate}[a)]
    \item Το εσωτερικό γινόμενο των διανυσμάτων $\vec{\alpha}, \vec{\beta}$ δίνεται από το τύπο
    \[ \vec{\alpha} \cdot \vec{\beta} = |\vec{\alpha}||\vec{\beta}| \sigma\upsilon\nu \left(\widehat{\vec{\alpha}, \vec{\beta}}\right) = 3 \cdot 4 \cdot \sigma\upsilon\nu \frac{\pi}{3} = 12 \cdot \frac{1}{2} = 6. \]
    Άρα $\vec{\alpha} \cdot \vec{\beta} = 6.$
    \item $\vec{\alpha}^2 = |\vec{\alpha}|^2 = 3^2 = 9$ και $\vec{\beta}^2 = |\vec{\beta}|^2 = 4^2 = 16.$
    \item $(3\vec{\alpha} - \vec{\beta}) \cdot (\vec{\alpha} - 3\vec{\beta}) = 3\vec{\alpha}^2 - \vec{\beta} \cdot \vec{\alpha} - 9\vec{\alpha} \cdot \vec{\beta} + 3\vec{\beta}^2 = 3\vec{\alpha}^2 - 10\vec{\alpha} \cdot \vec{\beta} + 3\vec{\beta}^2 = $
    \[ = 3 \cdot 9 - 10 \cdot 6 + 3 \cdot 16 = 15. \]
\end{enumerate}
\end{lysh}